% ============================================================
% Section 48: 一维晶格的总动量与声子准动量
% ============================================================
\newpage
\section{固体理论：一维晶格的总动量与声子准动量}
\label{sec:phonon-quasimomentum}

\begin{modelbox}[题目：晶格振动的总动量与声子准动量的含义]
考虑由 $N$ 个相同原子组成的一维晶格，晶格常数为 $a$，原子质量为 $m$，最近邻恢复力常数为 $\beta$。
\begin{enumerate}
    \item 求晶格振动的总动量；
    \item 讨论声子准动量 $\hbar q$ 的物理含义。
\end{enumerate}
\end{modelbox}

\begin{tipbox}[考点快速分析]
\begin{itemize}
    \item \textbf{总动量 vs 准动量}：$q\neq0$ 的格波总动量为零，但具有准动量 $\hbar q$
    \item \textbf{周期性边界条件}：$q=2\pi l/(Na)$，$l$ 为整数
    \item \textbf{等比级数求和}：$\sum_{n=0}^{N-1}e^{iqna}=0$（$q\neq0$）或 $N$（$q=0$）
    \item \textbf{准动量守恒}：源于晶格离散平移对称性，而非连续平移对称性
\end{itemize}
\end{tipbox}

\begin{derivationbox}[标准解题过程]
\textbf{Step 1: 格波解与单原子动量}

第 $n$ 个原子的运动方程为
\begin{equation}
m\ddot x_n = \beta(x_{n+1}+x_{n-1}-2x_n).
\end{equation}
设格波解 $x_n=A\,e^{i(qna-\omega t)}$，则原子动量
\begin{equation}
p_n = m\dot x_n = -i\omega mA\,e^{-i\omega t}e^{iqna}.
\end{equation}

\textbf{Step 2: 晶格总动量}

总动量
\begin{equation}
P = \sum_{n=0}^{N-1}p_n = -i\omega mA\,e^{-i\omega t}\sum_{n=0}^{N-1}e^{iqna}.
\label{eq:total-momentum}
\end{equation}

采用周期性边界条件（Born--von Karman），波矢量子化：
\begin{equation}
q = \frac{2\pi l}{Na}, \qquad l=0,\pm1,\pm2,\dots
\end{equation}

\textbf{Step 3: 等比级数求和}

式 \eqref{eq:total-momentum} 中的求和为等比级数：
\begin{equation}
\sum_{n=0}^{N-1}e^{iqna} = \frac{1-e^{iqNa}}{1-e^{iqa}} = \frac{1-e^{i2\pi l}}{1-e^{i2\pi l/N}}.
\end{equation}

\textbf{情况一：$q\neq0$（即 $l\neq0$）}

分子 $1-e^{i2\pi l}=0$，分母 $1-e^{i2\pi l/N}\neq0$（因 $l/N$ 非整数），故
\begin{equation}
\boxed{P=0 \quad (q\neq0).}
\end{equation}

\textit{物理意义}：$q\neq0$ 的格波模式描述的是原子间的\textbf{相对振动}，各原子动量相位不同，相互抵消，质心不动。

\textbf{情况二：$q=0$（即 $l=0$）}

分子分母同时为零，需取极限。直接代入 $q=0$：
\begin{equation}
\sum_{n=0}^{N-1}e^{i0} = N.
\end{equation}
于是
\begin{equation}
P = -i\omega mA\,e^{-i\omega t}\cdot N = (Nm)\cdot(-i\omega A\,e^{-i\omega t}) = M_{\text{total}}\,\dot x_{\text{cm}}.
\end{equation}
其中 $M_{\text{total}}=Nm$ 为整条链的总质量，$x_{\text{cm}}=A\,e^{-i\omega t}$ 为质心位移。
\begin{equation}
\boxed{P = M_{\text{total}}\,v_{\text{cm}} \quad (q=0).}
\end{equation}

\textit{物理意义}：$q=0$ 模式对应所有原子\textbf{同相运动}，即晶体的整体刚性平移，携带真实的物理动量。
\end{derivationbox}

\begin{formulabox}[结果汇总]
\begin{center}
\begin{tabular}{ccc}
\toprule
波矢 & 总动量 & 物理意义 \\
\midrule
$q=0$ & $P=M_{\text{total}}v_{\text{cm}}$ & 整体平移，携带\textbf{真实物理动量} \\
$q\neq0$ & $P=0$ & 内部相对振动，\textbf{不携带物理动量} \\
\bottomrule
\end{tabular}
\end{center}
\end{formulabox}

\begin{insightbox}[声子准动量 $\hbar q$ 的深层含义]
\textbf{1. 准动量不是物理动量}

$q\neq0$ 的声子模式总动量严格为零，但我们在声子--电子、声子--光子相互作用中却赋予它一个``动量''$\hbar q$。这是\textbf{准动量}，不是真实的机械动量。

\textbf{2. 准动量的来源——离散平移对称性}

在连续空间中，平移对称性导致动量守恒（Noether 定理）。在晶体中，只有\textbf{离散}平移对称性（平移 $a$ 晶格复原），因此守恒量不是连续动量，而是\textbf{在倒格子意义下的守恒}：
\begin{equation}
\sum \vec q_{\text{初}} = \sum \vec q_{\text{末}} + \vec G,
\end{equation}
其中 $\vec G$ 为任意倒格矢。这就是\textbf{准动量守恒}，允许相差一个倒格矢。

\textbf{3. 为什么需要引入准动量？}

在相互作用过程中（如电子-声子散射、中子非弹性散射），准动量守恒是\textbf{选择定则}的核心：
\begin{itemize}
    \item 中子吸收一个声子：$\vec k_{\text{中子}}' = \vec k_{\text{中子}} + \vec q + \vec G$
    \item 光子 Raman 散射：$\vec k_{\text{光子}}' = \vec k_{\text{光子}} \pm \vec q + \vec G$
\end{itemize}
如果没有准动量概念，就无法解释这些散射过程中波矢的选择规则。

\textbf{4. $q=0$ 的光学模是否有物理动量？}

$q=0$ 的光学模描述原胞内不同原子反相振动，质心仍不动，因此总动量也为零。只有 $q=0$ 的\textbf{声学模}（所有原子同相）才携带物理动量。

\textbf{5. 与德布罗意关系的类比}

虽然声子不携带真实动量，但 $\hbar q$ 在形式上满足与光子动量完全类似的守恒律。这提示我们：准动量是晶格周期性势场中波包的\textbf{晶体动量}（crystal momentum），类似于 Bloch 电子的动量 $\hbar k$。
\end{insightbox}

\begin{thinkbox}[考场速记：总动量与准动量的辨析模板]
遇到``讨论声子准动量与真实动量区别''类问题：
\begin{enumerate}
    \item \textbf{写总动量}：$P=\sum_n p_n$，代入格波解，化为等比级数
    \item \textbf{求和分情况}：$q\neq0$ 时级数和为零；$q=0$ 时和为 $N$
    \item \textbf{结论}：$q=0$ 有物理动量（整体平移），$q\neq0$ 动量为零
    \item \textbf{准动量解释}：源于离散平移对称性，守恒律允许相差倒格矢 $\vec G$
    \item \textbf{实验意义}：中子散射、Raman 散射的选择定则由准动量守恒决定
\end{enumerate}
\textit{核心口诀}：真实动量看质心，准动量看波矢；声子振动质心不动，故无真实动量，但有准动量。
\end{thinkbox}
